% Приложение Б — Промпт-шаблоны экспериментов
\setmonofont{Courier New}
\chapter{ПРОМПТ-ШАБЛОНЫ ЭКСПЕРИМЕНТОВ}

В данном приложении приведены полные тексты промпт-шаблонов, использованных во всех трёх экспериментах. Для каждого эксперимента указаны системный промпт (system), пользовательский промпт (user) и все варианты условий. Тексты воспроизведены дословно из экспериментального кода.

% -------------------------------------------------------
\section{Эксперимент 1: Промпты задачи Линды}
\label{sec:appB_linda}
% -------------------------------------------------------

\subsection{Системный промпт}

Системный промпт задаёт формат ответа и подавляет пояснения:

\begin{prompt}
You are solving a probabilistic choice task. Strictly follow the user's response format and do not add explanations. Return exactly two lines: <<CHOICE: a>> or <<CHOICE: b>>, then <<CONFIDENCE: N>>. Do not use parentheses or extra symbols around the choice letter.
\end{prompt}

\subsection{Пользовательский промпт}

Пользовательский промпт предъявляет виньетку и два варианта ответа. Ниже приведён шаблон с плейсхолдерами:

\begin{prompt}
Read the following description and answer the question.
{vignette}
Which is more probable?
(a) {simple_option}
(b) {conjunction_option}
First, indicate your choice: just the letter (a) or (b).
Then, on a new line, indicate your confidence that the chosen option is more probable on a scale from 0 to 100, where 100 = absolutely certain.
Answer strictly in two lines.
Do not add explanations.
Do not use parentheses around a or b.
Do not use square brackets in the response itself.
Response format:
CHOICE: [a or b]
CONFIDENCE: [0--100]
\end{prompt}

Здесь \texttt{\{vignette\}} --- подставленный текст виньетки с конкретными демографическими значениями (имя, раса, возраст); \texttt{\{simple\_option\}} --- вариант T (профессия); \texttt{\{conjunction\_option\}} --- вариант T\_and\_F1 (в условии A) или T\_and\_F2 (в условии B).

\subsection{Пример инстанцированного промпта}

Подстановка для виньетки ``public defender'' с именем Wei, расой Asian, возрастом 38, условие A (correlated):

\begin{prompt}
Read the following description and answer the question.
Wei is an Asian 38-year-old public defender who takes on cases others refuse. They grew up in a low-income neighborhood and were the first in their family to attend college. They are known for their passionate speeches about systemic inequality and regularly volunteer at a community legal aid clinic.
Which is more probable?
(a) public defender
(b) public defender who volunteers at a civil rights organization
First, indicate your choice: just the letter (a) or (b).
Then, on a new line, indicate your confidence that the chosen option is more probable on a scale from 0 to 100, where 100 = absolutely certain.
Answer strictly in two lines.
Do not add explanations.
Do not use parentheses around a or b.
Do not use square brackets in the response itself.
Response format:
CHOICE: [a or b]
CONFIDENCE: [0--100]
\end{prompt}

Все модели запрашивались с temperature~=~0. Каждая виньетка тестировалась в двух независимых сессиях (условия A и B) без переноса контекста.

% -------------------------------------------------------
\section{Эксперимент 2: Промпты задачи выбора карточек Уэйсона}
\label{sec:appB_wason}
% -------------------------------------------------------

Эксперимент использовал три варианта системного промпта и единый пользовательский шаблон. Все условия --- zero-shot.

\subsection{Вариант 1: Original (Assisted Falsification)}

Системный промпт задаёт ролевое прайминирование и явный фрейминг фальсификации:

\begin{prompt}
You are a logic expert analyzing conditional rules.
Your task is to identify which cards must be turned over to check if a rule is violated.

CRITICAL INSTRUCTIONS:
- Respond ONLY with the selected cards as a comma-separated list
- NO explanations, NO reasoning, NO JSON, NO other text
- Example: A, 4
- Another example: K, 7, A
\end{prompt}

\subsection{Вариант 2: Neutral}

Минимальный промпт без персоны и без фрейминга фальсификации:

\begin{prompt}
Respond ONLY with card labels as a comma-separated list.
No explanations.
\end{prompt}

\subsection{Вариант 3: Chain-of-Thought (CoT)}

Промпт с пошаговым рассуждением и фиксированным форматом финального ответа:

\begin{prompt}
You are a logic expert analyzing conditional rules.
Think step by step before answering.
After reasoning, output your final answer on the last line as:
ANSWER: card1, card2
\end{prompt}

\subsection{Пользовательский промпт}

Единый шаблон для всех трёх вариантов системного промпта:

\begin{prompt}
Rule: "{rule_text}"
Cards: {ordered_cards}
Which cards must you turn over?
Your answer:
\end{prompt}

Здесь \texttt{\{rule\_text\}} --- текст условного правила; \texttt{\{ordered\_cards\}} --- список из четырёх карточек в одной из пяти пермутаций. Все модели запрашивались с temperature~=~1{,}0 для обеспечения вариабельности ответов между перестановками.

% -------------------------------------------------------
\section{Эксперимент 3: Промпты задачи открытия правила 2-4-6}
\label{sec:appB_246}
% -------------------------------------------------------

Эксперимент использовал два варианта системного промпта (Baseline и Adaptive) и единый протокол диалога.

\subsection{Baseline: системный промпт}

\begin{prompt}
You have a hidden rule for number triples (x,y,z).
The triple (2,4,6) satisfies this rule.

Your goal: discover the rule through testing.

YOU MUST respond in a fenced python block on EVERY turn, using this EXACT content:
H: lambda x,y,z: <hypothesis>
T: [x, y, z]

STRICT RULES:
1. You MUST provide both H and T on every single turn
2. Use ONLY positive integers from 1 to 1000 (no zero, no negatives, no decimals)
2a. NEVER use 1001 or larger - any value above 1000 is invalid
2b. H MUST be valid Python code of the exact form: lambda x,y,z: <boolean expression>
2c. Use Python operators only: ==, !=, <, <=, >, >=, +, -, *, /, %, and, or, not
2d. NEVER use plain "=" inside H. Equality must be written as "=="
2e. If H is not valid Python lambda syntax, your answer is INVALID
3. NO explanations, commentary, or reasoning in your response
4. NO stopping until maximum turns reached
5. If you respond without a python block, your response is INVALID
6. If you respond without "T: [...]", your response is INVALID

Example response:
```python
H: lambda x,y,z: x < y < z
T: [5, 10, 15]
```
\end{prompt}

\subsection{Adaptive: системный промпт}

Идентичен Baseline по техническим ограничениям и формату, но дополнен блоком метакогнитивных указаний:

\begin{prompt}
[... идентичная базовая часть ...]

CRITICAL TESTING APPROACH:
- Your hypothesis should explain ALL observed results (both True and False)
- If you get a False result, your hypothesis MUST change to exclude that case
- If you get an unexpected True, your hypothesis MUST expand to include it
- Prefer tests that could prove your current hypothesis wrong, not just confirm it
- Repeatedly proposing triples your current hypothesis already expects to be True is weak testing
- Use boundary cases and contrastive cases that distinguish your current hypothesis from nearby alternatives
- Consider mathematical properties: sums, products, differences, divisibility, remainders (% operator)
- The rule is specific - not all triples satisfy it
- Pay attention to what makes some triples return True and others False
\end{prompt}

\subsection{Протокол диалога}

Диалог начинается с сообщения пользователя, подтверждающего порождающую тройку:

\begin{prompt}
User: True
\end{prompt}

На каждом последующем ходу пользователь предоставляет бинарную обратную связь:

\begin{prompt}
User: (x, y, z) -> True
\end{prompt}

или

\begin{prompt}
User: (x, y, z) -> False
\end{prompt}

Модель отвечает на каждом ходу в формате:

\begin{prompt}
```python
H: lambda x, y, z: <hypothesis>
T: [x, y, z]
```
\end{prompt}

Испытание завершается по достижении \texttt{max\_turns} или при $\text{IOU}(H_t, R) > 0{,}95$ --- в зависимости от того, что наступит раньше. Модели reasoning-типа (например, \texttt{gpt-5.2} reasoning) дополнительно используют отдельный reasoning-токен, не входящий в основной промпт.
